\documentclass[11pt,a4paper]{article}
\usepackage[margin=2.6cm]{geometry}
\usepackage{amsmath,amssymb}
\usepackage{booktabs}
\usepackage{microtype}
\usepackage[hidelinks]{hyperref}
\usepackage{setspace}
\usepackage{titlesec}
\usepackage{caption}
\captionsetup{font=small,labelfont=bf,justification=raggedright,singlelinecheck=false}
\titleformat{\section}{\normalfont\bfseries}{\thesection}{0.6em}{}
\titleformat{\subsection}{\normalfont\itshape}{\thesubsection}{0.6em}{}
\setlength{\parskip}{0.4em}
\setlength{\parindent}{0pt}
\onehalfspacing

\title{\normalsize\bfseries Retail Demand for a Counter-Inflation Currency:\\
A Simulation Under Explicit Assumptions}
\author{\normalsize Category One Limited}
\date{\normalsize Model version 3 \quad\textbullet\quad Draft}

\begin{document}
\maketitle

\begin{abstract}
\noindent\small
This study simulates retail adoption of a monetary instrument offering
protection against loss of purchasing power. It asks what share of household
savings would move into such an instrument over ten years, given stated
assumptions about how households learn of it, gain access to it, and come to
trust it. The instrument's protective property is taken as an input and is
not tested. Seventeen behavioural parameters with no direct measurement were
sampled simultaneously from continuous ranges across 8{,}000 draws, split
between a developed-market and a high-inflation specification. Adoption
occurred in 99.92 per cent of developed-market samples and in all
high-inflation samples. Median ten-year outcomes were 13.2 and 12.5 per cent
of household savings balances respectively, with fifth percentiles of 1.1 and
3.9 per cent. The binding constraint differs by market: trust in developed
economies, awareness where inflation is high. No dataset exists on
general-public demand for an instrument of this kind, and the study does not
constitute one.
\end{abstract}

\section{Scope}

This is a simulation of retail adoption. It does not predict adoption; it
shows what follows from its assumptions. It does not validate the
instrument: full purchasing-power protection enters as a parameter
$\pi = 1$, taken as given per the system's specification, and every result
is conditional on that holding. It does not measure demand; no dataset on
general-public demand for an inflation-protected savings instrument exists,
and every available source either samples existing cryptocurrency users or
measures a different product. It does not model commercial or regulatory
feasibility, with the exception noted in Section~7.

The central result is not a point estimate. It is that adoption occurs
across essentially the entire sampled parameter space, including at its
pessimistic extremes.

\section{Model}

An agent-based discrete-time model over monthly steps $t = 0, \dots, T$ with
$T = 120$, run over populations of $n$ synthetic households per market.

Each household $i$ carries a savings balance $b_i$ drawn lognormally, a
current holding $h_i \in \{L, S, O\}$ denoting local currency, US-dollar
stablecoin, or other store of value, an annual savings-drawdown fraction
$d_i$, a trust threshold $\theta_i$, and binary states for awareness and
account access.

\subsection{Marginal gain}

The annual benefit from adoption depends on current holding, not on
location:
\begin{equation}
g_i =
\pi \cdot
\begin{cases}
\iota_L & h_i = L\\[2pt]
\iota_U + \rho\sigma & h_i = S\\[2pt]
\tfrac{1}{2}\iota_U + \rho\sigma & h_i = O
\end{cases}
\end{equation}
where $\iota_L$ is local inflation, $\iota_U$ US inflation, $\rho$ the
perceived annual probability of a dollar stress event and $\sigma$ its
severity.

This asymmetry is the model's central economic mechanic. A household holding
Argentine pesos gains approximately $0.30$; one already holding a
dollar-denominated stablecoin gains approximately $0.03$, identically in
Buenos Aires and in London. Treating these as one population is the
principal source of overstatement in conventional total-addressable-market
estimates.

\subsection{Access}

Day-one access is restricted to households already holding a cryptocurrency
exchange account. In-app threshold-signature wallets combined with
decentralised-exchange liquidity render purchase a two-step operation for
that group, requiring no listing decision by any venue. Households without
an account must first open one and complete identity verification.

Account penetration grows toward a ceiling:
\begin{equation}
a(t) = \min\!\left(a_0 (1+\gamma)^{t/12},\; \bar{a}\right)
\end{equation}
with $a_0$ current penetration, $\gamma$ annual growth and $\bar{a}$ the
long-run ceiling. Household $i$ has access at $t$ if $u_i < a(t)$, where
$u_i \sim U(0,1)$ is drawn once.

\subsection{Salience}

Protection-motivated switching is rapid when a threat is vivid and slow
otherwise. Salience combines a baseline, decaying endogenous events, and a
constant manufactured component:
\begin{equation}
s(t) = \min\!\left(s_0 + \sum_{k} m_k \, 2^{-(t - t_k)/\lambda}
\,\mathbf{1}[t \ge t_k] + c,\; 1\right)
\end{equation}
where $(t_k, m_k)$ are event times and magnitudes, $\lambda$ the half-life
in months, and $c$ the manufactured component.

\subsection{Trust}

The trust threshold is an absolute required excess return, reduced by
salience and by discrete validation events:
\begin{equation}
\theta_i(t) = \theta_i \left(1 - \tau s(t)\right)
\prod_{j : t \ge t_j} (1 - e_j)
\end{equation}
with $\tau$ the salience effect and $(t_j, e_j)$ validation events and their
effects. A household adopts only if $g_i - \phi(t) d_i > \theta_i(t)$, where
$\phi(t)$ is the prevailing round-trip conversion spread.

\subsection{Diffusion}

Awareness spreads at different rates within and across the account-holding
population. For household $i$ at step $t$, with $A(t)$ the adopted share of
balances and $A_a(t)$ the adopted share among account holders:
\begin{equation}
p_i(t) =
\begin{cases}
\left(\alpha + \beta_{\text{in}} A_a(t)\right)/12 & \text{$i$ holds an account}\\[3pt]
\left(\tfrac{1}{2}\alpha + \beta_{\text{ac}} A(t)\right)/12 & \text{otherwise}
\end{cases}
\end{equation}
where $\alpha$ is external reach and $\beta_{\text{in}}, \beta_{\text{ac}}$
are the within- and across-population word-of-mouth coefficients. The
asymmetry reflects that a listener already holding an account can act
immediately, whereas one without must first open an account.

\subsection{Allocation}

Allocation rises with experience rather than with the size of the benefit.
For a household having held for $m_i$ months:
\begin{equation}
w_i(t) = \min\!\left(w_0 + r \frac{m_i}{12},\; \bar{w},\;
1 - \min(d_i \kappa / 12,\, 0.9)\right)
\end{equation}
with $w_0$ the initial allocation, $r$ the annual ramp rate, $\bar{w}$ the
ceiling and $\kappa$ the working buffer in months.

This yields a testable implication: mean allocation across adopters remains
below the ceiling for as long as adoption is growing, because recent
entrants lower the average. In the baseline developed-market run, 48.4 per
cent of households have adopted by year ten while mean allocation stands at
39.8 per cent against a ceiling of 60 per cent.

\section{Calibration}

Diffusion speed is anchored on M-Pesa rather than on cryptocurrency
adoption. M-Pesa reached 43 per cent of Kenyan households within seventeen
months of its March 2007 launch and approximately 70 per cent within four
years. Contemporary accounts attribute this to its agent network, which grew
from 4{,}000 to 19{,}000 cash-in and cash-out outlets against roughly 850
bank branches nationally, rather than to advertising. The inference drawn is
that distribution infrastructure, not awareness, was the binding constraint.

The trust threshold is set low but non-zero. Cohen and Dupas found
insecticide-treated bed-net take-up near 75 per cent when free, falling to
approximately 20 per cent at a small positive price; a discontinuity at
zero rather than an elasticity. A no-premium protective instrument should
therefore face a low barrier. It should not face none: a free bed net has no
downside, whereas moving savings into a novel instrument carries perceived
principal risk.

Salience decay is anchored on the observed pattern in catastrophe insurance,
where take-up rises after a nearby event and decays over subsequent years.

\section{Data}

\begin{table}[h]
\small
\begin{tabular}{@{}ll@{}}
\toprule
Quantity & Source \\
\midrule
US household savings deposits & Federal Reserve Z.1, \texttt{BOGZ1FL193030205Q}, Q1 2026 \\
UK household deposits & IMF Financial Access Survey, \texttt{GBRFCLODCHXDC} \\
Japan household cash and deposits & Bank of Japan Flow of Funds, Q1 2026 \\
China household savings & People's Bank of China \\
India, Turkey, Argentina, Chile & IMF Financial Access Survey \\
Euro area household deposits & Derived from European Central Bank \\
Inflation rates & National consumer price indices \\
\bottomrule
\end{tabular}
\end{table}

Approximately 70 per cent of the savings-pool total is cited. The remainder
is derived as nominal GDP multiplied by a household-deposit ratio banded
against the two Financial Access Survey ratios retrieved, being India at
39.9 per cent and Turkey at approximately 19 per cent. Brazil, Indonesia and
Mexico do not report a household breakdown to the Survey, so complete
citation is not attainable from this source.

The denominator is household savings, not M2. M2 includes business deposits,
transaction balances and money market fund shares, none of which are
addressable; substituting it approximately doubles every figure.

\subsection{Account penetration}

Published estimates disagree materially, and the access parameter is
sensitive to which is adopted. For the United States: Security.org's 2026
consumer report gives 30 per cent of adults from 992 respondents; the
National Cryptocurrency Association with Harris Poll gives approximately 25
per cent; Motley Fool Money gives 22 per cent from 2{,}000 respondents; Pew
Research Center gives approximately 20 per cent having ever used, from
8{,}512 respondents; and the Federal Reserve's Survey of Household Economics
and Decisionmaking, from a nationally representative sample of approximately
13{,}000 adults surveyed in October 2025, gives 10 per cent having used or
held within the past year, of which 7 per cent held as an investment.

The Federal Reserve figure is the most methodologically rigorous and is
three times lower than the commercial surveys. The discrepancy is partly
definitional, since the Survey measures activity within the past year
whereas the commercial surveys measure current ownership, and partly a
matter of sampling frames.

What this model requires is neither quantity exactly. Access requires a
funded exchange account, which persists after a holder sells, and is
therefore closer to having ever owned than to past-year activity. The
baseline adopts the upper end on that reasoning; the parameter is swept, and
this disagreement is the strongest argument for doing so.

\section{Uncertainty}

Behavioural parameters have no direct measurement, because no instrument of
this kind has been offered to this population. Each is classified as cited,
being traceable to a published figure; derived, being computed from cited
figures with the calculation stated; anchored, being set against a
documented reference class; or unmeasured, having a plausible range only and
being swept rather than fixed. No parameter falls outside these categories.
The full classification accompanies the code.

\subsection{Global sensitivity}

Seventeen parameters were sampled simultaneously from continuous ranges
across 4{,}000 draws per market rather than varied individually.
One-at-a-time sweeping holds remaining parameters at their defaults and
systematically overstates the leading parameter's influence; an earlier
version of this work reported a leverage ratio inflated by that method.

\begin{table}[h]
\small
\begin{tabular}{@{}lrr@{}}
\toprule
Parameter & Developed & High inflation \\
\midrule
Trust threshold & $-0.511$ & $-0.132$ \\
External reach & $+0.287$ & $+0.408$ \\
Allocation ceiling & $+0.273$ & $+0.379$ \\
Account ceiling & $+0.255$ & $+0.349$ \\
Salience effect on trust & $+0.221$ & $+0.140$ \\
Account growth & $+0.211$ & $+0.320$ \\
Allocation ramp rate & $+0.197$ & $+0.283$ \\
Salience baseline & $+0.147$ & $+0.023$ \\
Manufactured salience & $+0.146$ & $+0.014$ \\
Word of mouth, within population & $+0.104$ & $+0.143$ \\
Round-trip spread & $-0.104$ & $-0.041$ \\
Churn & $-0.097$ & $-0.143$ \\
Savings drawdown & $-0.078$ & $-0.013$ \\
Initial allocation & $+0.071$ & $+0.105$ \\
Salience half-life & $+0.063$ & $+0.056$ \\
Adoption hazard & $+0.044$ & $+0.060$ \\
Word of mouth, across populations & $+0.023$ & $+0.028$ \\
\bottomrule
\end{tabular}
\caption{Spearman rank correlation with ten-year share of household savings
balances. 4{,}000 draws per market.}
\end{table}

No parameter dominates. The binding constraint differs by market: trust in
developed economies, at approximately twice the influence of any other
parameter, and awareness where inflation is high. Manufactured salience
registers in developed markets and is negligible in high-inflation ones,
consistent with a campaign being unable to add vividness to a threat already
experienced directly.

\subsection{Outcome distribution}

\begin{table}[h]
\small
\begin{tabular}{@{}lrr@{}}
\toprule
& Developed & High inflation \\
\midrule
Minimum & 0.04\% & 0.99\% \\
5th percentile & 1.12\% & 3.87\% \\
25th percentile & 6.13\% & 8.09\% \\
Median & 13.21\% & 12.54\% \\
75th percentile & 21.63\% & 18.31\% \\
95th percentile & 37.82\% & 30.63\% \\
Maximum & 73.87\% & 58.04\% \\
\midrule
Samples with adoption & 99.92\% & 100.00\% \\
Samples self-sustaining at year ten & 92.0\% & 98.8\% \\
\bottomrule
\end{tabular}
\caption{Distribution of ten-year outcomes as a share of household savings
balances, across 4{,}000 draws per market.}
\end{table}

Across 8{,}000 samples spanning seventeen simultaneously varied unmeasured
parameters, adoption occurred in all but three.

\section{Reproducibility}

A run is fully determined by its configuration file and seed. No global
random number generator, no wall-clock time and no environment variable
enters the model. Each run writes a manifest recording the configuration
hash, seed, commit hash, library versions and a checksum of every output.
Reference fixtures are committed, and the test suite fails where a differing
environment produces differing figures. Code, configurations and the
parameter register are published so results may be reproduced, modified or
contested.

\section{Limitations}

Categorical access barriers are treated as scenario boundaries. Markets
whose capital controls prohibit the instrument are excluded from the base
case and reported separately, in the manner of an epidemiological study
excluding a population to which an intervention cannot lawfully be offered.
Jurisdiction-specific compliance questions are out of scope as
implementation-dependent.

China is reported separately and is never included in an aggregate total. At
approximately \$22.5 trillion it is the largest single pool and would
dominate any aggregate containing it.

Early-year figures carry the widest uncertainty. The initial segment of a
diffusion curve is governed almost entirely by unmeasured parameters; small
changes move the first year by multiples while barely affecting the tenth.

Terminal figures depend heavily on two unmeasured ceilings, being allocation
and account penetration, each swept but neither measured. Account
penetration estimates vary threefold across published sources; a reader
preferring the Federal Reserve figure should read the lower tail of the
distribution rather than the median.

No competitive or regulatory response is modelled.

\section{Corrections to earlier versions}

Four findings reported during development were subsequently identified as
artefacts rather than results.

A reported source-pool asymmetry of between 4:1 and 47:1 between
local-currency holders and stablecoin holders arose from an error in the
friction calculation, which annualised conversion cost by transaction count
rather than by the portion of the balance withdrawn, rendering the
stablecoin cohort's net advantage negative. Corrected, both cohorts adopt.

A reported allocation asymmetry of 3.75 was the ratio of two configuration
values set by hand, being an input restated rather than a model output.
Across a range moving that ratio fivefold, the ten-year result varied by 17
per cent.

A reported awareness threshold, below which adoption fails to become
self-sustaining, is a property of the diffusion form rather than a feature
of the instrument. Its location moves with the word-of-mouth coefficient and
disappears at high values.

A reported parameter leverage ratio of 13.7 was produced by one-at-a-time
sweeping. Under simultaneous sampling no parameter dominates.

\section{Measurements that would improve this work}

A landing-page conversion test in two target markets with real product
terms, which would measure the developed-market binding constraint directly.
A general-population survey screened for savings balance, being the only
route to demand data uncontaminated by cryptocurrency-user sampling.
Household penetration data in place of transaction-volume growth for
diffusion calibration. Market-maker spread indications, once a venue exists.

\vspace{1em}
\hrule
\vspace{0.6em}
{\footnotesize
Nothing in this document constitutes an offer to sell or a solicitation of
an offer to buy any security or financial instrument, nor investment,
financial, legal or tax advice. Outputs are conditional projections derived
from stated assumptions, many of them unmeasured. They are not forecasts.
\par}

\end{document}
